\section{Reductions and Undecidability}

\subsection{The Concept of Reductions}

\begin{remark}
    \textbf{Core Idea:} Reductions are the primary tool in computability and complexity theory to compare the relative difficulty of algorithmic problems.
    If a problem $A$ can be solved using an algorithm (subroutine) for a problem $B$, we say that \textbf{$A$ reduces to $B$} ($A \le B$). This means:
    $$
        \text{Problem } B \text{ is at least as hard as problem } A.
    $$
\end{remark}

\begin{important}
    \textbf{The Two Ways to Use a Reduction ($A \le B$):}
    \begin{itemize}
        \item \textbf{Positive Direction (Algorithm Design / Decidability):} If problem $B$ is known to be decidable, then problem $A$ is also decidable:
        $$
            A \le B \;\land\; B \in \mathcal{L}_{\text{rec}} \implies A \in \mathcal{L}_{\text{rec}}
        $$
        \item \textbf{Negative Direction (Impossibility Proofs / Undecidability):} By contraposition, if problem $A$ is known to be undecidable, then problem $B$ must also be undecidable:
        $$
            A \le B \;\land\; A \notin \mathcal{L}_{\text{rec}} \implies B \notin \mathcal{L}_{\text{rec}}
        $$
        This negative direction is the main method used to establish that newly encountered computational problems are undecidable.
    \end{itemize}
\end{important}

\subsection{General Reductions ($R$-Reductions)}

\begin{definition}
    \textbf{$R$-Reduction (Turing / Oracle Reduction):} Let $L_1 \subseteq \Sigma_1^*$ and $L_2 \subseteq \Sigma_2^*$. We say that $L_1$ is \textbf{$R$-reducible} to $L_2$ (denoted $L_1 \le_R L_2$) if there exists an algorithm (or Turing machine with an oracle for $L_2$) that decides membership in $L_1$, using queries to an imaginary subroutine for $L_2$ as a black box.
\end{definition}

\begin{theorem}
    \textbf{Properties of $R$-Reductions:}
    \begin{itemize}
        \item \textbf{Reflexivity and Transitivity:} $L \le_R L$, and if $L_1 \le_R L_2$ and $L_2 \le_R L_3$, then $L_1 \le_R L_3$.
        \item \textbf{Decidability Transfer:} If $L_1 \le_R L_2$ and $L_2 \in \mathcal{L}_{\text{rec}}$, then $L_1 \in \mathcal{L}_{\text{rec}}$.
        \item \textbf{Complement Invariance:} $L \le_R L^c$ always holds (simply invert the yes/no answer of the subroutine).
    \end{itemize}
\end{theorem}

\subsection{Many-One Reductions (EE-Reductions)}

\begin{definition}
    \textbf{EE-Reduction (Eingabe-Eingabe / Many-One Reduction):} A language $L_1 \subseteq \Sigma_1^*$ is called \textbf{EE-reducible} to $L_2 \subseteq \Sigma_2^*$ (written $L_1 \le_{EE} L_2$ or $L_1 \le_m L_2$) if there exists a \textbf{computable total function}
    $$
        f : \Sigma_1^* \to \Sigma_2^*
    $$
    such that for every word $w \in \Sigma_1^*$:
    $$
        w \in L_1 \iff f(w) \in L_2
    $$
    The function $f$ is computed by a halting deterministic Turing machine $M_f$ (a transducer).
\end{definition}

\begin{remark}
    \textbf{Intuition behind the Name:} It is called an \textit{Eingabe-Eingabe-Reduktion} (input-to-input reduction) because the reduction function $f$ directly maps an input $w$ of problem $L_1$ into a corresponding input $f(w)$ of problem $L_2$, without any post-processing of the decision result:
    $$
        \text{Decider for } L_1(w) = \text{Decider for } L_2(f(w))
    $$
    Every EE-reduction is a special case of an $R$-reduction: $L_1 \le_{EE} L_2 \implies L_1 \le_R L_2$.
\end{remark}

\begin{theorem}
    \textbf{Conservation Theorems for EE-Reductions:} Let $L_1 \le_{EE} L_2$ with reduction function $f$.
    \begin{enumerate}
        \item \textbf{Decidability:}
        \begin{itemize}
            \item If $L_2 \in \mathcal{L}_{\text{rec}}$, then $L_1 \in \mathcal{L}_{\text{rec}}$.
            \item If $L_1 \notin \mathcal{L}_{\text{rec}}$, then $L_2 \notin \mathcal{L}_{\text{rec}}$.
        \end{itemize}
        \item \textbf{Turing-Recognizability:}
        \begin{itemize}
            \item If $L_2 \in \mathcal{L}_{\text{RE}}$, then $L_1 \in \mathcal{L}_{\text{RE}}$.
            \item If $L_1 \notin \mathcal{L}_{\text{RE}}$, then $L_2 \notin \mathcal{L}_{\text{RE}}$.
        \end{itemize}
        \item \textbf{Complements:}
        $$
            L_1 \le_{EE} L_2 \iff L_1^c \le_{EE} L_2^c \quad (\text{using the exact same function } f)
        $$
    \end{enumerate}
\end{theorem}

\begin{remark}
    \textbf{Difference between $\le_R$ and $\le_{EE}$ on Recognizability:}
    While $L_1 \le_R L_2$ preserves decidability, it does \textbf{not} necessarily preserve Turing-recognizability (for instance, $L_{\text{diag}} \le_R L_{\text{diag}}^c$ holds by flipping the result, but $L_{\text{diag}}^c \in \mathcal{L}_{\text{RE}}$ while $L_{\text{diag}} \notin \mathcal{L}_{\text{RE}}$). In contrast, an EE-reduction \textbf{preserves both} decidability and Turing-recognizability.
\end{remark}

\subsection{Canonical Undecidable Problems and Classic Reductions}

\begin{definition}
    \textbf{The Halting Problem ($H$ and $H_\varepsilon$):}
    \begin{itemize}
        \item \textbf{General Halting Problem:}
        $$
            H = \{ \langle M \rangle \# w \mid M \text{ halts on input } w \}
        $$
        \item \textbf{Empty-Tape Halting Problem (Special Halting Problem):}
        $$
            H_\varepsilon = \{ \langle M \rangle \mid M \text{ halts on the empty input } \varepsilon \}
        $$
    \end{itemize}
\end{definition}

\begin{example}
    \textbf{Reduction $L_{\text{univ}} \le_{EE} H$:}
    Given an instance $\langle M \rangle \# w$ of $L_{\text{univ}}$, construct a Turing machine $M'$ that simulates $M$ on $w$. If $M$ accepts $w$, $M'$ halts. If $M$ rejects $w$, $M'$ enters an intentional infinite loop.
    Then:
    $$
        \langle M \rangle \# w \in L_{\text{univ}} \iff M \text{ accepts } w \iff M' \text{ halts} \iff \langle M' \rangle \# w \in H
    $$
    Since $L_{\text{univ}}$ is undecidable, it follows that $H$ is \textbf{undecidable}.
\end{example}

\begin{example}
    \textbf{Reduction $H \le_{EE} H_\varepsilon$:}
    Given an instance $\langle M \rangle \# w$ of $H$, define a reduction function $f(\langle M \rangle \# w) = \langle M_w \rangle$, where $M_w$ is a newly constructed Turing machine that:
    \begin{enumerate}
        \item Ignores its own input tape.
        \item Writes the fixed word $w$ onto its tape.
        \item Simulates $M$ on $w$.
    \end{enumerate}
    When started on the empty word $\varepsilon$, $M_w$ writes $w$ and simulates $M(w)$. Thus:
    $$
        M_w \text{ halts on } \varepsilon \iff M \text{ halts on } w
    $$
    Therefore, $\langle M \rangle \# w \in H \iff \langle M_w \rangle \in H_\varepsilon$.
    Because $H$ is undecidable, $H_\varepsilon$ is also \textbf{undecidable}.
\end{example}

\subsection{The Emptiness Problem and Non-Recognizability}

\begin{definition}
    \textbf{The Emptiness Problem ($L_\emptyset$):} The language of all Turing machines that accept no words:
    $$
        L_\emptyset = \{ \langle M \rangle \mid L(M) = \emptyset \}
    $$
    Its complement is $L_{\ne \emptyset} = \{ \langle M \rangle \mid L(M) \neq \emptyset \}$.
\end{definition}

\begin{theorem}
    \textbf{Status of $L_\emptyset$ and $L_{\ne \emptyset}$:}
    \begin{itemize}
        \item $L_{\ne \emptyset}$ is \textbf{Turing-recognizable} ($L_{\ne \emptyset} \in \mathcal{L}_{\text{RE}}$), via dovetailing simulation across all possible words.
        \item $L_{\ne \emptyset}$ is \textbf{undecidable} ($H_\varepsilon \le_{EE} L_{\ne \emptyset}$).
        \item $L_\emptyset$ is \textbf{not Turing-recognizable} ($L_\emptyset \notin \mathcal{L}_{\text{RE}}$).
    \end{itemize}
\end{theorem}

\subsection{Rice's Theorem}

\begin{definition}
    \textbf{Semantic Property of Turing Machines:} A property of Turing machines is called a \textbf{semantic property} (or language property) if it depends solely on the recognized language $L(M)$, not on the specific encoding or implementation details of $M$:
    $$
        L(M_1) = L(M_2) \implies (\langle M_1 \rangle \in L_{\mathcal{P}} \iff \langle M_2 \rangle \in L_{\mathcal{P}})
    $$
    A property $\mathcal{P} \subseteq \mathcal{L}_{\text{RE}}$ is called \textbf{non-trivial} if there is at least one recognizable language that has property $\mathcal{P}$ and at least one that does not:
    $$
        \mathcal{P} \neq \emptyset \quad \text{and} \quad \mathcal{P} \neq \mathcal{L}_{\text{RE}}
    $$
\end{definition}

\begin{theorem}
    \textbf{Rice's Theorem (Satz von Rice):} Every non-trivial semantic property of Turing-recognizable languages is \textbf{undecidable}.
    That is, if $\emptyset \subsetneq \mathcal{P} \subsetneq \mathcal{L}_{\text{RE}}$, then:
    $$
        L_{\mathcal{P}} = \{ \langle M \rangle \mid L(M) \in \mathcal{P} \} \notin \mathcal{L}_{\text{rec}}
    $$
\end{theorem}

\begin{remark}
    \textbf{Power of Rice's Theorem:} Rice's theorem immediately proves the undecidability of virtually every semantic question about program behavior without needing separate manual reductions:
    \begin{itemize}
        \item Is $L(M) = \emptyset$? (Undecidable)
        \item Is $L(M)$ finite? (Undecidable)
        \item Is $L(M)$ regular? (Undecidable)
        \item Does $M$ accept the word $01$? (Undecidable)
    \end{itemize}
    \textbf{Warning (Syntactic properties are NOT covered):} Rice's theorem only applies to semantic properties (languages). Syntactic properties (e.g., ,,Does $M$ have more than 5 states?`` or ,,Does $M$ ever execute transition $\delta(q_1, 0)$ on its transition diagram?``) are structural features of the code and are typically easily decidable.
\end{remark}
